\chapter{模型校核、性能评价方法与工况设计}

\section{基准工况自检与模型校核}

进行多工况仿真分析之前，要先针对着陆控制动力学模型进行基准工况自检和模型校核工作，目的是确保后续结果具备可比性与可信性。本文把平坦地形、正常模式、无初始侧偏、无额外航向误差选作基准工况。其参数设定如下：初始离地高度为120 ft，初始速度是72 kts，初始下降率为-250 fpm，跑道距离为2200 ft，下滑角为3.0°，目标接地点设定在跑道阈值后120 ft处，同时开启约束保护逻辑。

依据仿真结果可知，在基准工况时飞行器能够沿着预定的下滑轨迹实现稳定下降，其离地高度随着时间持续减小，速度维持在安全范围之内，俯仰角和滚转角的变化较小，这表明该模型在纵向控制和侧向控制方面整体比较协调。当进入着陆末段之后，系统能够及时从进近阶段转换至拉平阶段，并且在接地区域优先执行稳定接地逻辑，没有出现明显的轨迹突变或者状态发散情况。已有结果显示，在该工况下仿真时长约为17.15 s，末段离地高度约为4.60 ft，末段速度约为71.80 kts，末段俯仰角约为-1.28°，末段滚转角约为-2.28°。这些数据表明模型在标准条件下运行稳定，能够满足后续复杂工况分析的需求。

模型校核主要有三个方向：其一为轨迹合理性，即查看飞行器能否沿着预定下滑路径持续靠近跑道；其二为状态边界合理性，即检验速度、下沉率、姿态角、侧偏量是否处于合理区间；其三为阶段切换合理性，即核查系统是否能够在适宜的高度和距离进入拉平阶段，并在近地区域维持以稳定接地为主的控制策略。

\section{评价指标与判据建立}

为系统评价不同工况下固定翼飞行器着陆控制系统的性能，本文建立了评价指标模式，该模式包括安全性、稳定性、精度和交互有效性这四个方面。安全性指标包括飞行速度、最大下沉率、俯仰角与滚转角。稳定性指标涵盖离地高度变化是否连续、轨迹是否平滑、阶段切换是否合理。精度指标有侧偏量、航向误差、目标接地点附近位置的控制情况。交互有效性指标包括提示信息、告警触发、约束保护是否及时有效。

依据仿真平台的各项参数，本文确定了以下基本判据。在飞行速度处于62至88节范围，并且最大下沉率不超过850英尺每分钟，同时滚转角绝对值不超过20度、俯仰角保持在负12度至16度之间的时候，判定系统符合基本飞行安全要求。当飞行器进入拉平阶段之后，如果侧偏量不超过28英尺、航向误差不超过6度、滚转角不超过10度，判定拉平阶段横向对准状态良好。当进入接地提交阶段后，侧偏量不超过45英尺、航向误差不超过10度、滚转角不超过12度，便认为系统满足稳定接地要求。

\begin{table}[H]
\centering
\caption{着陆性能评价指标与判据}
\begin{tabular}{|l|l|l|}
\hline
\textbf{指标类别} & \textbf{指标名称} & \textbf{判据} \\
\hline
安全性 & 飞行速度 & 62~88 kts \\
\hline
安全性 & 最大下沉率 & 不大于850 fpm \\
\hline
安全性 & 滚转角 & |roll|≤20° \\
\hline
安全性 & 俯仰角 & -12°~16° \\
\hline
精度 & 拉平阶段侧偏量 & ≤28 ft \\
\hline
精度 & 拉平阶段航向误差 & ≤6° \\
\hline
精度 & 接地提交侧偏量 & ≤45 ft \\
\hline
精度 & 接地提交航向误差 & ≤10° \\
\hline
稳定性 & 阶段切换 & 平滑、无明显突变 \\
\hline
交互有效性 & 告警与约束触发 & 及时、准确 \\
\hline
\end{tabular}
\end{table}

\section{工况矩阵设计}

为剖析不同条件下模型的响应规律，本文依据已有平台参数和仿真结果，设计了包括"地形条件+控制模式+边界扰动"的工况矩阵。在地形条件方面，挑选了平坦地形、起伏地形、山脊地形、谷地地形、阶跃地形五种具有代表性的环境；在控制模式上，设定了正常模式和保守模式两种模式；针对边界扰动，可以增加初始侧偏、航向误差或者较大下沉率等因素，从而建立更为复杂的工况场景。

\begin{table}[H]
\centering
\caption{仿真工况矩阵设计}
\begin{tabular}{|l|l|l|l|l|l|}
\hline
\textbf{工况编号} & \textbf{地形类型} & \textbf{控制模式} & \textbf{初始侧偏} & \textbf{初始下降率} & \textbf{分析目标} \\
\hline
C1 & 平坦地形 & normal & 0 ft & -250 fpm & 基准工况校核 \\
\hline
C2 & 起伏地形 & normal & 0 ft & -250 fpm & 连续高程变化影响 \\
\hline
C3 & 山脊地形 & normal & 0 ft & -250 fpm & 近地抬升影响 \\
\hline
C4 & 谷地地形 & normal & 0 ft & -250 fpm & 谷地环境影响 \\
\hline
C5 & 阶跃地形 & normal & 0 ft & -250 fpm & 地形突变影响 \\
\hline
C6 & 平坦地形 & conservative & 0 ft & -250 fpm & 保守速度策略效果 \\
\hline
C7 & 起伏地形 & normal & 20 ft & -250 fpm & 横向偏差影响 \\
\hline
C8 & 山脊地形 & normal & 0 ft & -400 fpm & 较大下沉率影响 \\
\hline
\end{tabular}
\end{table}

\section{仿真计算与数据整理}

在工况矩阵确定完毕之后，本文运用统一的参数口径来进行仿真计算工作，并且借助CSV导出和曲线图绘制的方式对结果加以整理。仿真输出所涵盖的内容主要有时间、飞行距离、跑道相对距离、目标接地点距离、离地高度、海拔高度、速度、下沉率、俯仰角、滚转角、航向角、侧偏量、控制指令、飞行阶段、提示与告警信息等。

\begin{table}[H]
\centering
\caption{不同地形基础工况关键结果整理}
\begin{tabular}{|l|l|l|l|l|l|}
\hline
\textbf{地形工况} & \textbf{仿真时长/s} & \textbf{末段AGL/ft} & \textbf{末段速度/kts} & \textbf{末段俯仰角/(°)} & \textbf{末段滚转角/(°)} \\
\hline
平坦地形 & 17.15 & 4.60 & 71.80 & -1.28 & -2.28 \\
\hline
起伏地形 & 14.32 & 4.70 & 71.45 & -3.29 & -3.45 \\
\hline
山脊地形 & 14.03 & 4.74 & 70.44 & -2.52 & -2.96 \\
\hline
谷地地形 & 18.47 & 4.60 & 71.78 & -0.90 & 2.09 \\
\hline
阶跃地形 & 15.02 & 1.20 & 70.94 & -2.33 & -3.67 \\
\hline
\end{tabular}
\end{table}

由表可以看出，不同地形条件下系统输出存在明显差异。其中，起伏地形和山脊地形使飞行器姿态修正幅度更大，谷地地形使仿真持续时间相对较长，而阶跃地形对末段离地高度影响最突出。这些结果表明，本文所建立的模型和仿真平台能够较好地区分不同工况下的着陆响应差异。
